\section{Game Logic and Server Authority}
\label{sec:authority}

\begin{guidance}[Game logic and server authority (15 marks)]
Show: correct \textbf{move validation}; rejection of \textbf{illegal,
out-of-turn and out-of-state} moves with specific errors; reliable
\textbf{win and draw detection}; and a \textbf{client-untrusted} design. Most
importantly, give an argument that a \textbf{rejected request cannot change the
server state}, based on how your code is structured (e.g. every check runs
before the first write), and point to the tests that check it. The brief calls
this a ``proof''; make it a step-by-step argument a reader can verify against
your code.
\end{guidance}

\subsection{\ph{Heading about move validation}}
\ph{Which checks a move goes through, in which order, and where in the code
they live. A short listing of the validation code is useful here.}

\subsection{\ph{Heading about your argument that rejections change nothing}}
\ph{Step-by-step: where every error response is produced, and why no state is
written on any path that ends in an error. Cover framing errors, unknown
commands and every command handler. Then name the tests that check it.}

\subsection{\ph{Heading about turn control}}
\ph{How the server tracks and enforces whose turn it is, and how out-of-turn
requests are rejected.}

\subsection{\ph{Heading about win and draw detection}}
\ph{How you detect every win and the draw (e.g. a table of winning lines), in
which order you check them, and why that order matters. How the server refuses
moves once the game is over.}

\subsection{\ph{Heading about the client}}
\ph{Why the client contains no game rules, and what would go wrong if it did.}
